% Theorem numbers for Efimov_Large_Categories refer to arXiv:2405.12169v4.
	\chapter{Stable categories and localization sequences}
	\label[chapter]{chap:Localization_Sequences}

	In this chapter we collect the basic theory of localization sequences, both for small stable categories and for presentable stable categories, and introduce localizing invariants. Throughout, we denote by $\Cat_{\st}$ the category of small stable categories and exact functors, and by $\Cat_{\perf} \subseteq \Cat_{\st}$ the full subcategory of idempotent complete ones (this category is denoted $\Cat^{\perf}$ in \cite{BGT2013Universal, Efimov_Large_Categories}). For a stable category $B$ we write $\hom_B(x,y) \in \Sp$ for the mapping spectra. A full subcategory of a stable category is called \emph{stable} if it contains $0$ and is closed under fibers and cofibers, and \emph{thick} if it is moreover closed under retracts.

	\section{Verdier sequences}
	\label[section]{sec:Verdier_Sequences}

	% Source: written from the literature (Nikolaus--Scholze, Hermitian K-theory II)

	The basic construction of this chapter is the quotient of a small stable category $B$ by a stable subcategory $A$, obtained by forcing the objects of $A$ to become zero. In the setting of stable categories this is simply a localization in the sense of Dwyer--Kan, and the following theorem shows that it behaves just like Verdier's classical quotient of triangulated categories: it is again stable, it has the expected universal property, and its mapping spectra are computed by a colimit over $A$. We then introduce the corresponding notion of exact sequence, the Verdier sequences.

	\begin{theorem}[{\cite[Theorem~I.3.3]{NikolausScholze2018Cyclic}}]
		\label[theorem]{thm:Verdier_Quotient}
		Let $B$ be a small stable category and $A \subseteq B$ a full stable subcategory. Let $W_A$ be the class of morphisms of $B$ whose cofiber lies in $A$, and let $q\colon B \to B/A := B[W_A^{-1}]$ be the localization. Then:
		\begin{enumerate}
			\item The category $B/A$ is stable and $q$ is exact.
			\item For every stable category $E$, restriction along $q$ induces an equivalence between $\Fun^{\ex}(B/A,E)$ and the full subcategory of $\Fun^{\ex}(B,E)$ spanned by the functors that vanish on $A$.
			\item For $x,y \in B$, the category $A_{/y}$ is filtered and
			\[
				\Map_{B/A}(qx,qy) \simeq \colim_{(a \to y) \in A_{/y}} \Map_B\big(x,\cofib(a \to y)\big).
			\]
		\end{enumerate}
		The category $B/A$ is called the \emph{Verdier quotient} of $B$ by $A$.
	\end{theorem}

	Note that $q$ is essentially surjective, being a localization. Moreover, part (2) says that $B \to B/A$ is the cofiber of $A \hookrightarrow B$ in $\Cat_{\st}$. Indeed, for a stable category $E$, exact functors out of the pushout $0 \sqcup_A B$ in $\Cat_{\st}$ are given by
	\[
		\Fun^{\ex}(B,E) \times_{\Fun^{\ex}(A,E)} \{0\},
	\]
	the category of exact functors $F\colon B \to E$ together with an equivalence $F|_A \simeq 0$. Since $0$ is a zero object of $\Fun^{\ex}(A,E)$, the space of such equivalences is either empty or contractible, so this is the full subcategory of $\Fun^{\ex}(B,E)$ spanned by the functors vanishing on $A$.

	\begin{lemma}
		\label[lemma]{lem:Kernel_Verdier_Quotient}
		In the situation of \Cref{thm:Verdier_Quotient}, the kernel of $q\colon B \to B/A$ is the thick closure of $A$ in $B$, i.e.\ the full subcategory of retracts of objects of $A$. In particular, $B/A \simeq B/\bar{A}$, where $\bar{A}$ is the thick closure of $A$.
	\end{lemma}
	\prooffor{lem:Kernel_Verdier_Quotient}
	\begin{proof}
		Since $q$ is exact, its kernel is a thick subcategory of $B$, and it contains $A$ since $0 \to a$ lies in $W_A$ for $a \in A$. Conversely, let $x \in B$ with $q(x) \simeq 0$. Then $\id_{qx}$ is nullhomotopic, so by \Cref{thm:Verdier_Quotient}(3) (applied with $y = x$) there exists $a \to x$ with $a \in A$ such that the canonical map $x \to c := \cofib(a \to x)$ is nullhomotopic. Since $a \simeq \fib(x \to c)$ and the fiber of a nullhomotopic map $x \to c$ is $x \oplus \Omega c$, it follows that $x$ is a retract of $a$. This shows that the kernel of $q$ consists of retracts of objects of $A$, and every retract of an object of $A$ lies in the thick subcategory $\ker(q)$. The last statement follows from the universal property \Cref{thm:Verdier_Quotient}(2), since an exact functor vanishes on $A$ if and only if it vanishes on $\bar{A}$.
	\end{proof}

	\begin{remark}
		\label[remark]{rem:Nullhomotopies_Property}
		Let $F\colon B \to C$ be a functor into a pointed category $C$. Since the constant functor $0$ with value a zero object is a terminal object of $\Fun(B,C)$, the space of natural transformations $F \to 0$ is contractible. Hence the space of equivalences $F \simeq 0$ is either empty or contractible, and it is non-empty if and only if $F$ sends every object to a zero object. So for a composable pair $A \xrightarrow{i} B \xrightarrow{p} C$ of functors between pointed categories, a nullhomotopy $pi \simeq 0$ is unique if it exists, and its existence is a \emph{property} of the pair $(i,p)$. We will therefore write ``$pi \simeq 0$'' as a condition throughout, without specifying a nullhomotopy.
	\end{remark}

	\begin{definition}
		\label[definition]{def:Verdier_Sequence}
		A sequence $A \xrightarrow{i} B \xrightarrow{p} C$ in $\Cat_{\st}$ with $pi \simeq 0$ is called a \emph{Verdier sequence} if it is both a fiber sequence and a cofiber sequence in $\Cat_{\st}$. Explicitly, this means that $i$ is fully faithful with essential image $\ker(p)$, and that the induced functor $B/A \to C$ is an equivalence. In this case $p$ is called a \emph{Verdier projection}.
	\end{definition}

	By \Cref{lem:Kernel_Verdier_Quotient}, for every full stable subcategory $A \subseteq B$ the sequence $\bar{A} \hookrightarrow B \to B/A$ is a Verdier sequence. See \cite[Lemma~A.1.8 and Proposition~A.1.9]{Hermitian_II} for a detailed discussion of Verdier sequences.

	\section{Idempotent completion and short exact sequences}
	\label[section]{sec:Karoubi_Sequences}

	% Source: written from the literature (Lurie HA, Efimov, Hermitian K-theory II); Mathew, Lecture 1 (MPIM Bonn, 09.09.2024)

	Every small stable category $A$ admits an idempotent completion $A \hookrightarrow A^{\idem}$, which is again stable \cite[Corollary~1.1.3.7]{lurie2016HA}; it is fully faithful and every object of $A^{\idem}$ is a retract of an object of $A$, and $A \mapsto A^{\idem}$ is left adjoint to the inclusion $\Cat_{\perf} \subseteq \Cat_{\st}$ (see \Cref{sec:Idempotent_Completeness}). In particular, cofibers in $\Cat_{\perf}$ are computed as $(B/A)^{\idem}$.

	\begin{definition}
		\label[definition]{def:Short_Exact_Sequence}
		A sequence $A \xrightarrow{i} B \xrightarrow{p} C$ in $\Cat_{\perf}$ with $pi \simeq 0$ is called a \emph{short exact sequence} (or \emph{Karoubi sequence}) if it is both a fiber sequence and a cofiber sequence in $\Cat_{\perf}$. Explicitly, this means that $i$ is fully faithful with essential image $\ker(p)$, and that the induced functor $(B/A)^{\idem} \to C$ is an equivalence. In this case $p$ is called a \emph{Karoubi projection}.
	\end{definition}

	We follow \cite{Efimov_Large_Categories} in calling these short exact sequences; the term Karoubi sequence is used in \cite{Hermitian_II}.

	\begin{lemma}
		\label[lemma]{lem:Characterization_Short_Exact}
		A sequence $A \xrightarrow{i} B \xrightarrow{p} C$ in $\Cat_{\perf}$ with $pi \simeq 0$ is short exact if and only if $i$ is fully faithful with essential image $\ker(p)$, the induced functor $B/A \to C$ is fully faithful, and every object of $C$ is a retract of an object in the image of $p$.
	\end{lemma}
	\prooffor{lem:Characterization_Short_Exact}
	\begin{proof}
		The functor $\bar{p}\colon B/A \to C$ factors as $B/A \hookrightarrow (B/A)^{\idem} \to C$. If $\bar{p}$ is fully faithful, then so is $(B/A)^{\idem} \to C$, since mapping spaces between retracts of objects are retracts of mapping spaces between those objects. If moreover every object of $C$ is a retract of an object $p(b) = \bar{p}(qb)$, then $(B/A)^{\idem} \to C$ is essentially surjective since $(B/A)^{\idem}$ is idempotent complete. Conversely, if $(B/A)^{\idem} \to C$ is an equivalence, then $\bar{p}$ is fully faithful and every object of $C$ is a retract of an object of $B/A$, and hence of an object $p(b)$, as $q$ is essentially surjective.
	\end{proof}

	\begin{remark}
		\label[remark]{rem:Verdier_Vs_Karoubi}
		Every Verdier sequence between idempotent complete categories is short exact. The converse fails whenever $A \subseteq B$ is a thick subcategory such that the Verdier quotient $B/A$ is not idempotent complete: then $A \to B \to (B/A)^{\idem}$ is short exact but not a Verdier sequence.
	\end{remark}

	% Source: Mathew, Lecture 1 (MPIM Bonn, 09.09.2024)
	\begin{remark}
		\label[remark]{rem:Mapping_Spectra_Quotient}
		Let $A \subseteq B$ be a full stable subcategory with quotient functor $q\colon B \to B/A$. Applying \Cref{thm:Verdier_Quotient}(3) to all shifts of $x$, and using that the colimit is filtered and hence commutes with cofibers, we obtain the following formula for the mapping spectra:
		\begin{align*}
			\hom_{B/A}(qx,qy) &\simeq \colim_{(a \to y) \in A_{/y}} \hom_B\big(x,\cofib(a \to y)\big) \\
			&\simeq \cofib\Big(\colim_{(a \to y) \in A_{/y}} \hom_B(x,a) \to \hom_B(x,y)\Big).
		\end{align*}
		Since $B/A \to (B/A)^{\idem}$ is fully faithful, the same formula computes mapping spectra in the idempotent complete quotient.
	\end{remark}

	\begin{remark}
		\label[remark]{rem:Lifting_Objects_Quotient}
		In a short exact sequence $A \to B \xrightarrow{p} C$, the functor $p$ need not be essentially surjective: in general, objects of $C$ are only retracts of objects in the image of $p$. It is possible to say precisely which objects lift. By Thomason's classification of dense subcategories \cite{Thomason_Classification}, the essential image of the dense subcategory $B/A \subseteq C$ consists precisely of those objects $c \in C$ whose class $[c] \in K_0(C)$ lies in the image of $K_0(B) \to K_0(C)$. For example, $c \oplus \Sigma c$ always lifts to $B$, since $[c \oplus \Sigma c] = [c] - [c] = 0$.
	\end{remark}

	\begin{example}[{Thomason--Trobaugh, \cite{Thomason_Trobaugh}}]
		\label[example]{ex:Thomason_Trobaugh_Sequence}
		Let $X$ be a qcqs scheme and $U \subseteq X$ a quasi-compact open subscheme with closed complement $Z$. Then
		\[
			\Perf_Z(X) \hookrightarrow \Perf(X) \to \Perf(U)
		\]
		is a short exact sequence, where $\Perf_Z(X)$ denotes the perfect complexes on $X$ whose restriction to $U$ vanishes. In general, the restriction functor $\Perf(X) \to \Perf(U)$ is not essentially surjective; by \Cref{rem:Lifting_Objects_Quotient}, a perfect complex on $U$ extends to $X$ if and only if its class in $K_0(U)$ extends.
	\end{example}

	\section{Localization sequences of presentable stable categories}
	\label[section]{sec:Localization_Sequences_PrL}

	% Source: written from the literature (Lurie HTT, Efimov, Thomason--Neeman); Mathew, Lecture 1 (MPIM Bonn, 09.09.2024); Talk 1 (Marc Hoyois, 13.10.2026)

	We now turn to large categories. Recall that $\PrL_{\st}$ denotes the category of presentable stable categories and colimit-preserving functors. A colimit-preserving functor $F\colon D \to E$ between presentable categories has a right adjoint $F^R$ (\Cref{thm:Adjoint_Functor_Theorem}); recall that $F$ is called \emph{strongly continuous} if $F^R$ again preserves colimits. We denote by $\Pr^{\mathrm{LL}}_{\st} \subseteq \PrL_{\st}$ the subcategory with the same objects and strongly continuous functors.

	\begin{definition}
		\label[definition]{def:Localization_Sequence_PrL}
		A sequence $K \xrightarrow{i} D \xrightarrow{p} E$ in $\PrL_{\st}$ with $pi \simeq 0$ is called a \emph{localization sequence} if $i$ is fully faithful with essential image $\ker(p)$ and $p$ is a \emph{localization}, i.e.\ its right adjoint $p^R$ is fully faithful. It is called \emph{strongly continuous} if moreover $i$ and $p$ are strongly continuous.
	\end{definition}

	For a full subcategory $K \subseteq D$ we write $K^{\perp} = \{x \in D \mid \Map_D(k,x) \simeq * \text{ for all } k \in K\}$ for its right orthogonal.

	\begin{proposition}
		\label[proposition]{prop:Localization_Sequences_PrL}
		\leavevmode
		\begin{enumerate}
			\item Let $D$ be a presentable stable category and $K \subseteq D$ a full stable subcategory which is closed under colimits and presentable. Then the inclusion $i\colon K \hookrightarrow D$ has a right adjoint $i^R$, the full subcategory $K^{\perp} \subseteq D$ is presentable, its inclusion admits a left adjoint $L\colon D \to K^{\perp}$ given by
			\[
				L(x) = \cofib(ii^R x \to x),
			\]
			and $K \xrightarrow{i} D \xrightarrow{L} K^{\perp}$ is a localization sequence. It is both a fiber sequence and a cofiber sequence in $\PrL_{\st}$.
			\item Conversely, let $K \xrightarrow{i} D \xrightarrow{p} E$ be a localization sequence. Then $p^R$ induces an equivalence $E \simeq K^{\perp}$ under which $p$ corresponds to $L$, and there is a natural cofiber sequence
			\[
				ii^R \to \id_D \to p^Rp.
			\]
		\end{enumerate}
	\end{proposition}
	\prooffor{prop:Localization_Sequences_PrL}
	\begin{proof}
		(1) The inclusion $i$ preserves colimits, so it has a right adjoint by \Cref{thm:Adjoint_Functor_Theorem}. Define $L(x) = \cofib(ii^Rx \to x)$, using the counit. For $k \in K$, the map $\hom_D(k,ii^Rx) \to \hom_D(k,x)$ is an equivalence since $i$ is fully faithful, so $\hom_D(k,L(x)) \simeq 0$; that is, $L(x) \in K^{\perp}$. For $y \in K^{\perp}$, applying $\hom_D(-,y)$ to the cofiber sequence $ii^Rx \to x \to L(x)$ shows that $\hom_D(L(x),y) \to \hom_D(x,y)$ is an equivalence, since $\hom_D(ii^Rx,y) \simeq 0$. Hence $L\colon D \to K^{\perp}$ is left adjoint to the inclusion $K^{\perp} \subseteq D$, with unit $x \to L(x)$. The functor $L$ is accessible since $i^R$ is (being a right adjoint between presentable categories), so $K^{\perp}$ is presentable \cite[Section~5.5.4]{lurie2009HTT}. Since its right adjoint is fully faithful, $L$ is a localization, and $L(x) \simeq 0$ if and only if $ii^Rx \to x$ is an equivalence, if and only if $x \in K$. This shows that $K \to D \to K^{\perp}$ is a localization sequence. It is a fiber sequence in $\PrL_{\st}$ since limits in $\PrL$ are computed in $\widehat{\Cat}$. For the cofiber sequence, let $F\colon D \to E$ be a colimit-preserving functor to a presentable stable category $E$. Since $L$ is a localization, restriction along $L$ identifies $\Fun^{\colim}(K^{\perp},E)$ with the full subcategory of $\Fun^{\colim}(D,E)$ of functors inverting the maps $x \to x'$ with $L(x) \simeq L(x')$. Since $L$ is exact, these are precisely the maps whose fiber lies in $\ker(L) = K$, and an exact functor inverts all such maps if and only if it vanishes on $K$. This is precisely the universal property of the cofiber of $i$ in $\PrL_{\st}$.

		(2) Since $p^R$ is fully faithful and $\Map_D(k,p^Re) \simeq \Map_E(pk,e) \simeq *$ for $k \in K$, the functor $p^R$ lands in $K^{\perp}$. Conversely, let $x \in K^{\perp}$ and let $f \to x \to p^Rp(x)$ be the fiber sequence of the unit. As $pp^R \simeq \id$, we have $p(f) \simeq 0$, so $f \in K$. Then $\hom_D(f,x) \simeq 0$ and $\hom_D(f,p^Rpx) \simeq \hom_E(pf,px) \simeq 0$, so $\hom_D(f,f) \simeq 0$ and hence $f \simeq 0$. So $x \simeq p^Rp(x)$, which shows that $p^R$ induces an equivalence $E \simeq K^{\perp}$. Under this equivalence, $p$ is identified with the left adjoint $L$ of the inclusion $K^{\perp} \subseteq D$, and the unit $x \to p^Rp(x) = L(x)$ is the cofiber of $ii^Rx \to x$ by (1).
	\end{proof}

	\begin{lemma}
		\label[lemma]{lem:Strong_Continuity_Localization_Sequence}
		Let $K \xrightarrow{i} D \xrightarrow{p} E$ be a localization sequence. Then $i^R$ preserves colimits if and only if $p^R$ preserves colimits. In particular, the sequence is strongly continuous as soon as $i$ or $p$ is strongly continuous.
	\end{lemma}
	\prooffor{lem:Strong_Continuity_Localization_Sequence}
	\begin{proof}
		Since $i$ is fully faithful and preserves colimits, $i^R$ preserves colimits if and only if $ii^R$ does: indeed, for a diagram $(x_j)$ the map $i(\colim_j i^Rx_j) \simeq \colim_j ii^Rx_j \to ii^R(\colim_j x_j)$ is the image under $i$ of the comparison map for $i^R$. Similarly, $p^R$ preserves colimits if and only if $p^Rp$ does: if $p^Rp$ preserves colimits and $(e_j)$ is a diagram in $E$, then $\colim_j e_j \simeq p(\colim_j p^Re_j)$, so $p^R(\colim_j e_j) \simeq p^Rp(\colim_j p^Re_j) \simeq \colim_j p^Rpp^Re_j \simeq \colim_j p^Re_j$. Finally, by \Cref{prop:Localization_Sequences_PrL}(2) there is a cofiber sequence of exact functors $ii^R \to \id \to p^Rp$, so $ii^R$ preserves colimits if and only if $p^Rp$ does.
	\end{proof}

	Recall that a presentable stable category $E$ is \emph{compactly generated} if $E \simeq \Ind(E^{\omega})$, equivalently if there is a set $G$ of compact objects of $E$ such that an object $e \in E$ with $\hom_E(g,e) \simeq 0$ for all $g \in G$ is zero.

	\begin{lemma}
		\label[lemma]{lem:Compact_Generators_Thick_Closure}
		Let $E$ be a presentable stable category and $G \subseteq E^{\omega}$ a set of compact objects such that $\hom_E(g,e) \simeq 0$ for all $g \in G$ implies $e \simeq 0$. Then $E^{\omega}$ is the thick closure of $G$, and $E \simeq \Ind(E^{\omega})$.
	\end{lemma}
	\prooffor{lem:Compact_Generators_Thick_Closure}
	\begin{proof}
		Let $T \subseteq E^{\omega}$ be the thick closure of $G$; it is essentially small. The functor $\Ind(T) \to E$ extending the inclusion is fully faithful since $T$ consists of compact objects, and its essential image is closed under colimits since $T$ is closed under finite colimits \cite[Section~5.3.5]{lurie2009HTT}. Its right adjoint is conservative: if $\hom_E(t,e) \simeq 0$ for all $t \in T$ then $e \simeq 0$ by assumption. Hence $\Ind(T) \simeq E$. Since $T$ is idempotent complete, the compact objects of $\Ind(T)$ are precisely the objects of $T$ (\Cref{thm:Accessible_Recollections}(1)), so $E^{\omega} = T$.
	\end{proof}

	\begin{lemma}
		\label[lemma]{lem:Compact_Objects_Localization_Sequence}
		Let $K \xrightarrow{i} D \xrightarrow{p} E$ be a strongly continuous localization sequence. Then $i$ and $p$ preserve compact objects, and $K^{\omega} = K \cap D^{\omega}$. If $D$ is compactly generated, then so is $E$, and $E^{\omega}$ is the thick closure of $p(D^{\omega})$.
	\end{lemma}
	\prooffor{lem:Compact_Objects_Localization_Sequence}
	\begin{proof}
		A colimit-preserving functor whose right adjoint preserves filtered colimits preserves compact objects. If $k \in K$ is compact in $D$, then it is compact in $K$ since $i$ is fully faithful and preserves colimits. If $D$ is compactly generated and $e \in E$ satisfies $\hom_E(p(d),e) \simeq \hom_D(d,p^Re) \simeq 0$ for all $d \in D^{\omega}$, then $p^Re \simeq 0$ and hence $e \simeq pp^Re \simeq 0$. The claim follows from \Cref{lem:Compact_Generators_Thick_Closure}.
	\end{proof}

	\begin{definition}[{\cite[Definition~1.9]{Efimov_Large_Categories}}]
		\label[definition]{def:Homological_Epimorphism}
		A functor $f\colon A \to B$ in $\Cat_{\perf}$ is called a \emph{homological epimorphism} if $\Ind(f)\colon \Ind(A) \to \Ind(B)$ is a localization, i.e.\ its right adjoint is fully faithful.
	\end{definition}

	\begin{theorem}[Thomason--Neeman]
		\label[theorem]{thm:Thomason_Neeman}
		\leavevmode
		\begin{enumerate}
			\item If $A \xrightarrow{i} B \xrightarrow{p} C$ is a short exact sequence in $\Cat_{\perf}$, then $\Ind(A) \to \Ind(B) \to \Ind(C)$ is a strongly continuous localization sequence. In particular, every Karoubi projection is a homological epimorphism.
			\item Conversely, if $K \xrightarrow{i} D \xrightarrow{p} E$ is a strongly continuous localization sequence in which $K$ and $D$ are compactly generated, then $E$ is compactly generated and $K^{\omega} \to D^{\omega} \to E^{\omega}$ is a short exact sequence.
		\end{enumerate}
	\end{theorem}
	\prooffor{thm:Thomason_Neeman}
	\begin{proof}
		(1) The functor $\Ind(i)$ is fully faithful, as mapping spaces between ind-objects are computed by the formula of \Cref{constr:Ind_Large}. Its essential image $K \subseteq \Ind(B)$ is closed under colimits and presentable, so by \Cref{prop:Localization_Sequences_PrL}(1) we obtain a localization sequence $K \to \Ind(B) \xrightarrow{L} K^{\perp}$ which is a cofiber sequence in $\PrL_{\st}$. For every presentable stable category $E$, we have natural equivalences
		\[
			\begin{aligned}
				\Fun^{\colim}(K^{\perp},E) &\simeq \Fun^{\colim}_{K \mapsto 0}(\Ind(B),E) \simeq \Fun^{\ex}_{A \mapsto 0}(B,E) \\
				&\simeq \Fun^{\ex}(B/A,E) \simeq \Fun^{\ex}((B/A)^{\idem},E) \simeq \Fun^{\colim}(\Ind(C),E),
			\end{aligned}
		\]
		using \Cref{thm:Ind_Kappa_Universal_Property}(3), \Cref{thm:Verdier_Quotient}(2), the fact that $E$ is idempotent complete, and $(B/A)^{\idem} \simeq C$. Under these equivalences, precomposition with $L$ corresponds to precomposition with $\Ind(p)$. By the Yoneda lemma there is an equivalence $K^{\perp} \simeq \Ind(C)$ under which $L$ corresponds to $\Ind(p)$. Hence $\Ind(A) \to \Ind(B) \to \Ind(C)$ is a localization sequence. Since $\Ind(i)$ preserves compact objects, its right adjoint preserves filtered colimits, hence all colimits; so the sequence is strongly continuous by \Cref{lem:Strong_Continuity_Localization_Sequence}.

		(2) Put $A = K^{\omega}$ and $B = D^{\omega}$. By \Cref{lem:Compact_Objects_Localization_Sequence}, $i$ restricts to a fully faithful functor $A \to B$ whose image $A = K \cap B$ is a thick subcategory of $B$. By \Cref{lem:Kernel_Verdier_Quotient}, the sequence $A \to B \to (B/A)^{\idem}$ is short exact, so by (1) we obtain a localization sequence $\Ind(A) \to \Ind(B) \to \Ind((B/A)^{\idem})$. As $K \simeq \Ind(A)$ and $D \simeq \Ind(B)$, both $E$ and $\Ind((B/A)^{\idem})$ are identified with $K^{\perp}$ compatibly with the functors from $D$, by \Cref{prop:Localization_Sequences_PrL}(2). Hence $E \simeq \Ind((B/A)^{\idem})$ is compactly generated with $E^{\omega} \simeq (B/A)^{\idem}$, and the claim follows.
	\end{proof}

	Thus, by \Cref{thm:Thomason_Neeman}, $\Ind$ turns short exact sequences of small categories into strongly continuous localization sequences of compactly generated categories. Conversely, a strongly continuous localization sequence $K \to D \to E$ with $D$ compactly generated comes from a short exact sequence of compact objects if and only if $K$ is compactly generated; strong continuity alone does not guarantee this. In the language of \Cref{def:Homological_Epimorphism}: a functor $f\colon A \to B$ in $\Cat_{\perf}$ is a homological epimorphism if and only if $\ker(\Ind(f)) \to \Ind(A) \to \Ind(B)$ is a strongly continuous localization sequence, and it is a Karoubi projection if and only if moreover $\ker(\Ind(f))$ is compactly generated. In general, the kernel of a homological epimorphism need not be compactly generated, see \Cref{ex:Keller} below; it is, however, always a dualizable category (\Cref{chap:Compactly_Assembled_Dualizable}).

	% Source: Mathew, Lecture 1 (MPIM Bonn, 09.09.2024)
	\Cref{thm:Thomason_Neeman} gives a description of idempotent complete quotients in terms of large categories. Recall from \Cref{cor:Ind_As_Exact_Functors} that $\Ind(B) \simeq \Fun^{\ex}(B\catop,\Sp)$ for $B \in \Cat_{\perf}$.

	\begin{corollary}
		\label[corollary]{cor:Karoubi_Quotient_Via_Ind}
		Let $i\colon A \hookrightarrow B$ be a fully faithful functor in $\Cat_{\perf}$. Under the identification $\Ind(-) \simeq \Fun^{\ex}((-)\catop,\Sp)$, the right adjoint $\Ind(i)^R\colon \Ind(B) \to \Ind(A)$ is given by restriction along $i$, and
		\[
			(B/A)^{\idem} \simeq \ker\big(\Ind(i)^R\big)^{\omega},
		\]
		where $\ker(\Ind(i)^R) \subseteq \Fun^{\ex}(B\catop,\Sp)$ is the full subcategory of exact functors vanishing on $A$.
	\end{corollary}

	For the description of exact presheaves on $B/A$ by this vanishing condition, see \href{https://animarium.org/tag/1155/}{[SCT, Corollary~19.6.8]}.

	\prooffor{cor:Karoubi_Quotient_Via_Ind}
	\begin{proof}
		For $a \in A$ and $F \in \Ind(B)$ we have $\hom_{\Ind(A)}(a,\Ind(i)^RF) \simeq \hom_{\Ind(B)}(ia,F) \simeq F(ia)$ by the Yoneda lemma, so $\Ind(i)^R(F) \simeq F \circ i\catop$. Since $A$ is idempotent complete and $i$ is fully faithful, the essential image of $i$ is a thick subcategory, so $A \to B \to (B/A)^{\idem}$ is short exact by \Cref{lem:Kernel_Verdier_Quotient}. By \Cref{thm:Thomason_Neeman}(1), $\Ind(A) \to \Ind(B) \to \Ind((B/A)^{\idem})$ is a localization sequence, so by \Cref{prop:Localization_Sequences_PrL}(2) we have $\Ind((B/A)^{\idem}) \simeq \Ind(A)^{\perp}$. Finally, $F \in \Ind(A)^{\perp}$ if and only if $\hom(a,\Ind(i)^RF) \simeq \hom(ia,F) \simeq 0$ for all $a \in A$, if and only if $\Ind(i)^RF \simeq 0$. Taking compact objects gives the claim.
	\end{proof}

	% Source: Talk 1 (Marc Hoyois, 13.10.2026)
	\begin{remark}
		\label[remark]{rem:Ind_Fully_Faithful_Localization}
		The functor $\Ind$ induces an equivalence of categories
		\[
			\Ind\colon \Cat_{\perf} \rightleftarrows \PrL_{\omega,\st}\noloc (-)^{\omega},
		\]
		where $\PrL_{\omega,\st}$ denotes the category of compactly generated presentable stable categories and colimit-preserving functors that preserve compact objects (equivalently, strongly continuous functors), see \Cref{thm:Accessible_Recollections} and \Cref{lem:Compact_Generators_Thick_Closure}. For a functor $f\colon A \to B$ in $\Cat_{\perf}$:
		\begin{itemize}
			\item $f$ is fully faithful if and only if $\Ind(f)$ is fully faithful (by the formula for mapping spaces in \Cref{constr:Ind_Large}), if and only if $\Ind(f)$ is a monomorphism \cite[Proposition~1.13]{Efimov_Large_Categories};
			\item If $f$ is a Karoubi projection, then $\Ind(f)$ is a localization (\Cref{thm:Thomason_Neeman}), or equivalently $\Ind(f)$ is an epimorphism \cite[Proposition~1.13]{Efimov_Large_Categories}. The converse does not hold, see \Cref{ex:Keller}.
		\end{itemize}
	\end{remark}

	% Source: Talk 1 (Marc Hoyois, 13.10.2026)
	\begin{example}[{Keller, \cite{Keller_Smashing}}]
		\label[example]{ex:Keller}
		Let $A$ be a ring and let $0 \neq I \subseteq \mathrm{Jac}(A)$ be a Tor-unital ideal. Then $A/I \otimes^L_A A/I \simeq A/I$, so the restriction functor $\Mod_{A/I} \to \Mod_A$ is fully faithful and $f\colon \Perf_A \to \Perf_{A/I}$ is a homological epimorphism. The category
		\[
			\ker(\Mod_A \to \Mod_{A/I})
		\]
		contains no nonzero compact objects. Note that $\Mod_A \simeq \Ind(\Perf_A)$ and $\Mod_{A/I} \simeq \Ind(\Perf_{A/I})$. So the kernel of $f$ is zero, whereas the large kernel contains the nonzero module $I$, since $(A/I) \otimes_A^L I \simeq 0$. This gives a counterexample to the converse of the second statement of \Cref{rem:Ind_Fully_Faithful_Localization}: $\Ind(f)$ is a localization, but $f$ is not a Karoubi projection.

		For example, consider $A = k[t,t^{1/2},t^{1/4}, \dots]_{(t,t^{1/2}, t^{1/4}, \dots)}$, for $k$ a field. Let $\mm$ be the maximal ideal. Then the map $A \to A/\mm$ is an example.
	\end{example}

	\section{Semiorthogonal decompositions}
	\label[section]{sec:Semiorthogonal_Decompositions}

	% Source: written from the literature (Efimov, Section 2.11); Mathew, Lecture 1 (MPIM Bonn, 09.09.2024)

	The simplest localization sequences are the split ones. In this case the middle category $E$ is glued from two subcategories $C$ and $D$ between which there are no maps in one direction, and every object of $E$ is uniquely an extension of an object of $C$ by an object of $D$. Such a splitting is called a semiorthogonal decomposition. We follow Efimov's convention: in $E = \lra{C_1,\dots,C_n}$ there are no nonzero maps from $C_i$ to $C_j$ for $i > j$. We first treat the presentable setting, which is the one relevant for dualizable categories, and then explain the analogous notion for small stable categories and its relation to short exact sequences.

	\begin{definition}[{\cite[Definition~2.74]{Efimov_Large_Categories}}]
		\label[definition]{def:Semiorthogonal_Decomposition}
		A \emph{semiorthogonal decomposition} of a presentable stable category $E$ is a collection of full presentable stable subcategories $C_1,\dots,C_n \subseteq E$ such that
		\begin{enumerate}
			\item the inclusion functors $C_i \hookrightarrow E$ are strongly continuous;
			\item for $i > j$, $x \in C_i$ and $y \in C_j$ we have $\hom_E(x,y) \simeq 0$;
			\item the smallest full stable subcategory of $E$ containing $C_1,\dots,C_n$ is $E$ itself.
		\end{enumerate}
		In this case we write $E = \lra{C_1,\dots,C_n}$.
	\end{definition}

	\begin{proposition}
		\label[proposition]{prop:Two_Term_SOD}
		Let $E = \lra{C,D}$ be a semiorthogonal decomposition, with inclusions $i_C\colon C \hookrightarrow E$ and $i_D\colon D \hookrightarrow E$. Then:
		\begin{enumerate}
			\item Every $x \in E$ fits into a cofiber sequence $d \to x \to c$ with $d \in D$ and $c \in C$, which is unique up to unique equivalence and functorial in $x$.
			\item The functor $i_C$ admits a left adjoint $i_C^L$ and $i_D$ admits a right adjoint $i_D^R$, given by $x \mapsto c$ and $x \mapsto d$ respectively; there is a natural cofiber sequence $i_Di_D^R \to \id_E \to i_Ci_C^L$, and $i_C^Li_D \simeq 0$.
			\item The sequence $D \xrightarrow{i_D} E \xrightarrow{i_C^L} C$ is a strongly continuous localization sequence, split by $i_C$.
		\end{enumerate}
	\end{proposition}
	\prooffor{prop:Two_Term_SOD}
	\begin{proof}
		Note first that $\hom_E(d,c) \simeq 0$ for $d \in D$ and $c \in C$. Suppose $d \to x \to c$ is a cofiber sequence with $d \in D$ and $c \in C$. For $d' \in D$ we then have $\hom_E(d',d) \xrightarrow{\simeq} \hom_E(d',x)$, since $\hom_E(d',c) \simeq 0$; so $d \to x$ exhibits $d$ as the value at $x$ of a right adjoint to $i_D$. Dually, for $c' \in C$ we have $\hom_E(c,c') \xrightarrow{\simeq} \hom_E(x,c')$, so $x \to c$ exhibits $c$ as the value at $x$ of a left adjoint of $i_C$. In particular, such a sequence is unique up to unique equivalence and functorial in $x$, wherever it exists.

		Let $E' \subseteq E$ be the full subcategory of those $x$ that admit such a cofiber sequence. It contains $C$ (take $0 \to c \to c$) and $D$ (take $d \to d \to 0$) and is closed under shifts. It is closed under cofibers: a map $x \to x'$ in $E'$ induces, by the universal properties above, a map of cofiber sequences, and passing to cofibers gives a cofiber sequence $\cofib(d \to d') \to \cofib(x \to x') \to \cofib(c \to c')$ with first term in $D$ and last term in $C$, as these are stable subcategories. Hence $E'$ is a stable subcategory containing $C$ and $D$, so $E' = E$ by assumption (3). This proves (1), and by \cite[Proposition~5.2.4.2]{lurie2009HTT} also the existence of the adjoints in (2), together with the cofiber sequence $i_Di_D^R \to \id \to i_Ci_C^L$. Applying it to $x = i_D(d)$ gives $i_C^Li_D(d) \simeq 0$.

		For (3), $i_C^L$ is a localization since its right adjoint $i_C$ is fully faithful, and $i_C^L(x) \simeq 0$ if and only if $i_Di_D^R x \to x$ is an equivalence, if and only if $x$ lies in the essential image of $D$. The right adjoints $i_D^R$ and $i_C$ preserve colimits by assumption (1) and since $C \subseteq E$ is closed under colimits, so the sequence is strongly continuous. Finally, $i_C^L i_C \simeq \id_C$ as $i_C$ is fully faithful, so $i_C$ splits $i_C^L$.
	\end{proof}

	% Source: Talk 6 (Jacopo Mazzei, 17.11.2026); cf. Efimov, Remark 2.75
	Conversely, a semiorthogonal decomposition can be recognized from the adjoints of the inclusion functors. In particular, the projections onto the two pieces are themselves strongly continuous.

	\begin{proposition}[{\cite[Remark~2.75]{Efimov_Large_Categories}}]
		\label[proposition]{prop:SOD_Characterization}
		Let $i_1\colon C \to E$ and $i_2\colon D \to E$ be strongly continuous functors between presentable stable categories. Then $i_1$ and $i_2$ are fully faithful and exhibit $E = \lra{i_1(C), i_2(D)}$ as a semiorthogonal decomposition if and only if the following conditions hold:
		\begin{enumerate}
			\item $i_1$ admits a left adjoint $i_1^L$ and $i_2$ admits a right adjoint $i_2^R$;
			\item the counit $i_1^Li_1 \to \id_C$ and the unit $\id_D \to i_2^Ri_2$ are equivalences;
			\item $i_1^Li_2 \simeq 0$, or equivalently $i_2^Ri_1 \simeq 0$;
			\item the sequence $i_2i_2^R \to \id_E \to i_1i_1^L$, formed by the counit and the unit, is a cofiber sequence.
		\end{enumerate}
		In this case, the functors $i_1^L$ and $i_2^R$ are strongly continuous.
	\end{proposition}

	The composite in (4) admits a unique nullhomotopy, as $\Map(i_2i_2^R,i_1i_1^L) \simeq \Map(i_1^Li_2i_2^R,i_1^L) \simeq *$ by (3); so the condition in (4) makes sense. The equivalence of the two formulations of (3) follows from $\hom_D(d,i_2^Ri_1c) \simeq \hom_E(i_2d,i_1c) \simeq \hom_C(i_1^Li_2d,c)$.

	% Source: Talk 6 (Jacopo Mazzei, 17.11.2026); strong continuity of the projections added from Efimov, Remark 2.75
	\prooffor{prop:SOD_Characterization}
	\begin{proof}
		Suppose first that $i_1$ and $i_2$ are fully faithful and $E = \lra{i_1(C),i_2(D)}$. Then (1), (3) and (4) are \Cref{prop:Two_Term_SOD}(2), applied to the essential images, and (2) holds as $i_1$ and $i_2$ are fully faithful.

		Conversely, assume (1)--(4). By (2), the functors $i_1$ and $i_2$ are fully faithful. For $c \in C$ and $d \in D$ we have $\hom_E(i_2d,i_1c) \simeq \hom_C(i_1^Li_2d,c) \simeq 0$ by (3), so the semiorthogonality condition holds. For $e \in E$, the cofiber sequence $i_2i_2^Re \to e \to i_1i_1^Le$ of (4) shows that $e$ lies in the stable subcategory generated by the essential images of $i_1$ and $i_2$. Finally, the inclusions are strongly continuous by assumption, and their essential images are closed under colimits, as $i_1$ and $i_2$ preserve colimits.

		It remains to show that $i_1^L$ and $i_2^R$ are strongly continuous. The right adjoint of $i_1^L$ is $i_1$, which preserves colimits. The right adjoint of $i_1i_1^L$ is $i_1i_1^R$, which preserves colimits since $i_1$ is strongly continuous. Passing to right adjoints is an equivalence between the stable category of colimit-preserving endofunctors of $E$ and the opposite of the stable category of accessible limit-preserving endofunctors, so it carries the fiber sequence $i_2i_2^R \to \id_E \to i_1i_1^L$ to a cofiber sequence $i_1i_1^R \to \id_E \to (i_2i_2^R)^R$. Hence $(i_2i_2^R)^R$ preserves colimits. Since $i_2$ is fully faithful, we have
		\[
			\hom_D(i_2^Re,d) \simeq \hom_E(i_2i_2^Re,i_2d) \simeq \hom_E(e,(i_2i_2^R)^R i_2 d),
		\]
		so $(i_2i_2^R)^R \circ i_2$ is right adjoint to $i_2^R$ and preserves colimits.
	\end{proof}

	% Source: Mathew, Lecture 1 (MPIM Bonn, 09.09.2024)
	\begin{remark}
		\label[remark]{rem:SOD_Small}
		\Cref{def:Semiorthogonal_Decomposition} makes sense for any stable category $E$ and full stable subcategories $C_1,\dots,C_n \subseteq E$, omitting condition (1). No presentability or size assumption is needed. Parts (1) and (2) of \Cref{prop:Two_Term_SOD} hold in this setting with the same proof, and the sequence $D \xrightarrow{i_D} E \xrightarrow{i_C^L} C$ is a Verdier sequence: $i_C^L$ is a localization with kernel $D$, and it inverts precisely the maps whose cofiber lies in $D$. If $E$ is idempotent complete, then so are $C$ and $D$, being thick subcategories of $E$ (as kernels of the exact functors $i_D^R$ and $i_C^L$), and the sequence is short exact. (In the talk, the definition was given simultaneously for presentable stable categories, for dualizable categories and for small stable categories.)
	\end{remark}

	Conversely, short exact sequences admitting adjoints give rise to semiorthogonal decompositions.

	\begin{proposition}
		\label[proposition]{prop:SOD_From_Short_Exact}
		Let $A \xrightarrow{i} B \xrightarrow{p} C$ be a short exact sequence in $\Cat_{\perf}$. Then $i$ admits a right adjoint $i^R$ if and only if $p$ admits a right adjoint $p^R$. In this case:
		\begin{enumerate}
			\item The functor $p^R$ is fully faithful, with essential image the right orthogonal $A^{\perp} = \{x \in B \mid \hom_B(ia,x) \simeq 0 \text{ for all } a \in A\}$;
			\item There is a semiorthogonal decomposition $B = \lra{p^R(C),\, i(A)}$;
			\item The sequence $C \xrightarrow{p^R} B \xrightarrow{i^R} A$ is again short exact.
		\end{enumerate}
	\end{proposition}
	\prooffor{prop:SOD_From_Short_Exact}
	\begin{proof}
		Assume first that $i^R$ exists, and for $b \in B$ let $L(b) := \cofib(ii^Rb \to b)$. As in the proof of \Cref{prop:Localization_Sequences_PrL}(1), $L(b) \in A^{\perp}$ and $L\colon B \to A^{\perp}$ is left adjoint to the inclusion. Moreover, $L(b) \simeq 0$ if and only if $b \simeq ii^Rb$ lies in $A$, so $L$ inverts precisely the maps whose cofiber lies in $A$. Since $L$ is a reflective localization, it exhibits $A^{\perp}$ as the localization of $B$ at these maps \cite[Proposition~5.2.7.12]{lurie2009HTT}, so $B/A \simeq A^{\perp}$. As $A^{\perp}$ is closed under retracts in the idempotent complete category $B$, it is idempotent complete, so $C \simeq (B/A)^{\idem} \simeq A^{\perp}$, with $p$ corresponding to $L$. Hence $p$ admits the fully faithful right adjoint $A^{\perp} \hookrightarrow B$.

		Conversely, assume that $p^R$ exists. For $a \in A$ and $c \in C$ we have $\hom_B(ia,p^Rc) \simeq \hom_C(pia,c) \simeq 0$, so $p^R$ lands in $A^{\perp}$. Every $x \in A^{\perp}$ is local with respect to the maps $s\colon b' \to b$ with $\cofib(s) \in A$, i.e.\ $s^*\colon \hom_B(b,x) \to \hom_B(b',x)$ is an equivalence. Hence $\hom_{B/A}(qb,qx) \simeq \hom_B(b,x)$ for all $b \in B$ and $x \in A^{\perp}$ (for instance by \Cref{rem:Mapping_Spectra_Quotient}, applied to $B\catop$). Applying this to $x = p^Rc$, we find that the composite
		\[
			\hom_B(b,p^Rc) \to \hom_C(pb,pp^Rc) \to \hom_C(pb,c),
		\]
		whose second map is induced by the counit, is the adjunction equivalence. Its first map is an equivalence as well: since the functor $B/A \to C$ induced by $p$ is fully faithful, it is identified with the equivalence $\hom_B(b,p^Rc) \simeq \hom_{B/A}(qb,qp^Rc)$. Hence so is the second map, i.e.\ the counit $pp^Rc \to c$ induces an equivalence on $\hom_C(pb,-)$ for all $b \in B$. Since every object of $C$ is a retract of an object of the form $pb$, the counit is an equivalence, i.e.\ $p^R$ is fully faithful. For $b \in B$, let $f \to b \to p^Rpb$ be the fiber sequence of the unit. Then $p(f) \simeq 0$, so $f$ lies in $A$, and for $a \in A$ we get $\hom_B(ia,f) \simeq \hom_B(ia,b)$ since $\hom_B(ia,p^Rpb) \simeq 0$. Hence $b \mapsto f$ is right adjoint to $i$.

		It remains to prove (1)--(3). For (1), it remains to show that every $x \in A^{\perp}$ lies in the image of $p^R$. The fiber $f$ of the unit $x \to p^Rpx$ lies in $A$, so $\hom_B(f,x) \simeq 0$ and $\hom_B(f,p^Rpx) \simeq 0$, whence $\hom_B(f,f) \simeq 0$ and $f \simeq 0$. For (2), we have $\hom_B(ia,p^Rc) \simeq 0$, and every $b \in B$ fits into the cofiber sequence $ii^Rb \to b \to p^Rpb$. For (3), the kernel of $i^R$ is $A^{\perp} = p^R(C)$ by (1), and $i^R$ is a coreflective localization since its left adjoint $i$ is fully faithful; applying the first paragraph to opposite categories, $i^R$ exhibits $A$ as the Verdier quotient $B/p^R(C)$.
	\end{proof}

	\section{Calkin categories}
	\label[section]{sec:Calkin_Small}

	% Source: Mathew, Lecture 1 (MPIM Bonn, 09.09.2024)
	For $A \in \Cat_{\perf}$ and a regular cardinal $\kappa$, we write $\Ind(A)^{\kappa}$ for the full subcategory of $\kappa$-compact objects of $\Ind(A)$. It is essentially small and stable, and it contains $A = \Ind(A)^{\omega}$ as a thick subcategory.

	\begin{definition}[{\cite[Section~2.9]{Efimov_Large_Categories}}]
		\label[definition]{def:Calkin_Small}
		Let $A \in \Cat_{\perf}$ and let $\kappa$ be an uncountable regular cardinal. The \emph{$\kappa$-Calkin category} of $A$ is
		\[
			\Calk_{\kappa}(A) := \big(\Ind(A)^{\kappa}/A\big)^{\idem},
		\]
		so that there is a short exact sequence $A \hookrightarrow \Ind(A)^{\kappa} \to \Calk_{\kappa}(A)$.
	\end{definition}

	\begin{remark}
		\label[remark]{rem:Calkin_Mapping_Spectra}
		The mapping spectra in $\Calk_{\kappa}(A)$ can be computed using \Cref{rem:Mapping_Spectra_Quotient}. Let $x = \quotcolim_{i \in I} x_i$ and $y = \quotcolim_{j \in J} y_j$ be objects of $\Ind(A)^{\kappa}$, with $I$ and $J$ filtered. The functor $J \to A_{/y}$ sending $j$ to $y_j \to y$ is cofinal, and $\hom_{\Ind(A)}(x,y_j) \simeq \lim_i \hom_A(x_i,y_j)$. Hence
		\[
			\hom_{\Calk_{\kappa}(A)}(x,y) \simeq \cofib\Big(\colim_{j}\lim_{i}\hom_A(x_i,y_j) \to \lim_{i}\colim_{j}\hom_A(x_i,y_j)\Big).
		\]
		So the Calkin category measures the failure of the limit over $i$ to commute with the filtered colimit over $j$.

		To justify the cofinality assertion, fix $a \to y$ with $a \in A$. The comma category $(a \to y)/J$ is the category of elements of the functor
		\[
			j \mapsto \fib_{a \to y}\big(\Map_A(a,y_j) \to \Map_{\Ind(A)}(a,y)\big).
		\]
		Its groupoid completion is the colimit of this functor. Since $a$ is compact and filtered colimits of animae commute with finite limits, this colimit is the fiber of the identity of $\Map_{\Ind(A)}(a,y)$ over $a \to y$, hence is contractible. The cofinality criterion therefore applies.
	\end{remark}

	\begin{example}
		\label[example]{ex:Calkin_Perf}
		Let $R$ be an $\bbE_{\infty}$-ring and $A = \Perf_R$, so that $\Ind(A) \simeq \Mod_R$. For $M \in \Mod_R$ and a perfect $R$-module $P$ we have $\hom_R(M,P) \simeq \hom_R(M,R) \otimes_R P$, since $P$ is dualizable. Writing $M' \in \Mod_R$ as a filtered colimit of perfect modules, \Cref{rem:Mapping_Spectra_Quotient} gives
		\[
			\hom_{\Calk_{\kappa}(\Perf_R)}(M,M') \simeq \cofib\big(\hom_R(M,R) \otimes_R M' \to \hom_R(M,M')\big)
		\]
		for $M,M' \in \Mod_R^{\kappa}$. For instance, if $R = k$ is a field and $V$ is a $k$-vector space of dimension $< \kappa$, the map $V^{\vee} \otimes_k V \to \End_k(V)$ is injective with image the finite rank endomorphisms, so
		\[
			\End_{\Calk_{\kappa}(\Perf_k)}(V) \simeq \End_k(V)/\{\text{finite rank endomorphisms}\}.
		\]
		This is the algebraic analogue of the Calkin algebra $B(\Hh)/\Kk(\Hh)$ of bounded operators modulo compact operators on a Hilbert space, which explains the name.
	\end{example}

	\section{Localizing invariants}
	\label[section]{sec:Localizing_Invariants}

	% Source: Mathew, Lecture 1 (MPIM Bonn, 09.09.2024); Talk 1 (Marc Hoyois, 13.10.2026)
	We start with some recollections on the functor $K\colon \Cat_{\st} \to \Sp$.

	\begin{definition}
		\label[definition]{def:K_0}
		The \emph{Grothendieck group} $K_0(A)$ of a small stable category $A$ is the free abelian group on the equivalence classes $[x]$ of objects $x \in A$, modulo the relations $[x] = [x'] + [x'']$ for every cofiber sequence $x' \to x \to x''$ in $A$.
	\end{definition}

	Nonconnective algebraic K-theory is a spectrum-valued refinement $K\colon \Cat_{\st} \to \Sp$ of this construction, satisfying $\pi_0K(A) \cong K_0(A)$ for $A$ idempotent complete. Its fundamental property is that it is a localizing invariant in the following sense.

	\begin{definition}
		\label[definition]{def:Localizing_Invariant}
		Let $\Ee$ be a stable category. A functor $F\colon \Cat_{\perf} \to \Ee$ is called a \emph{localizing invariant} if $F(0) \simeq 0$ and $F$ sends short exact sequences (\Cref{def:Short_Exact_Sequence}) to fiber sequences. We extend $F$ to $\Cat_{\st}$ by setting $F(A) := F(A^{\idem})$. We denote by $\Fun_{\loc}(\Cat_{\st},\Ee)$ the full subcategory of localizing invariants.
	\end{definition}

	By \Cref{thm:Thomason_Neeman}, a sequence $A \to B \to C$ in $\Cat_{\perf}$ is short exact if and only if $\Ind(A) \to \Ind(B) \to \Ind(C)$ is a strongly continuous localization sequence in $\PrL_{\st}$.

	\begin{lemma}
		\label[lemma]{lem:Localizing_Invariants_Verdier}
		Let $\Ee$ be stable and let $F\colon \Cat_{\st} \to \Ee$ satisfy $F(0) \simeq 0$. Consider the following conditions:
		\begin{enumerate}
			\item For every Verdier sequence $A \hookrightarrow B \to B/A$ in the sense of \Cref{def:Verdier_Sequence}, $F$ takes this sequence to a fiber sequence.
			\item Each sequence $A \hookrightarrow B \to C$ with $A \to B$ fully faithful and inducing an equivalence $(B/A)^{\idem} \iso C^{\idem}$ is sent to a fiber sequence under $F$.
		\end{enumerate}
		Then (2) implies (1), and (2) holds if and only if $F$ satisfies (1) and inverts the maps $A \to A^{\idem}$. In this case, $F$ is (the extension of) a localizing invariant in the sense of \Cref{def:Localizing_Invariant}.
	\end{lemma}
	\prooffor{lem:Localizing_Invariants_Verdier}
	\begin{proof}
		The sequences in (1) are among those of (2). Suppose that $F$ satisfies (1) and inverts idempotent completions. For a sequence as in (2), the fully faithful functor $A^{\idem} \hookrightarrow B^{\idem}$ has thick image, so
		\[
			A^{\idem} \hookrightarrow B^{\idem} \to B^{\idem}/A^{\idem}
		\]
		is a Verdier sequence. Moreover, $(B^{\idem}/A^{\idem})^{\idem} \simeq (B/A)^{\idem} \simeq C^{\idem}$, by the universal property of the idempotent complete quotient. Applying (1) and idempotent invariance therefore gives (2). Conversely, (2) applied to $0 \to A \to A^{\idem}$ shows that $F(A) \to F(A^{\idem})$ is an isomorphism.
	\end{proof}

	\begin{remark}
		\label[remark]{rem:Verdier_Localizing_Not_Karoubi}
		Condition (1) of \Cref{lem:Localizing_Invariants_Verdier} alone does not imply that $F$ inverts $A \to A^{\idem}$. For instance, connective algebraic K-theory satisfies (1) (see the proof of \cite[Corollary~4.4.15]{Hermitian_II}), but $K_0(A) \to K_0(A^{\idem})$ is in general not surjective. This is why localizing invariants are defined on idempotent complete categories.
	\end{remark}

	% Source: Mathew, Lecture 1 (MPIM Bonn, 09.09.2024)
	\begin{remark}
		\label[remark]{rem:Finitary_Localizing_Invariants}
		One frequently additionally requires that $\Ee$ is presentable and that $F$ preserves filtered colimits (or at least $\kappa$-filtered colimits for some regular cardinal $\kappa$). Algebraic K-theory and $\THH$ are examples of localizing invariants preserving filtered colimits. Topological cyclic homology $\mathrm{TC}$ is localizing, but does not in general preserve filtered colimits.
	\end{remark}

	\begin{example}
		\label[example]{ex:Localizing_Additive}
		Let $B = \lra{C,D}$ be a semiorthogonal decomposition in $\Cat_{\perf}$ (\Cref{rem:SOD_Small}). Then every localizing invariant $F$ sends the short exact sequence $D \to B \to C$ to a fiber sequence $F(D) \to F(B) \to F(C)$, which is split by $F(i_C)$. Hence
		\[
			F(B) \simeq F(C) \oplus F(D).
		\]
		Functors $F\colon \Cat_{\perf} \to \Ee$ with this property are called \emph{additive}.
	\end{example}

	\begin{example}
		\label[example]{ex:Localizing_Calkin}
		Let $A \in \Cat_{\perf}$ and let $F$ be a localizing invariant. The short exact sequence $A \to \Ind(A)^{\omega_1} \to \Calk_{\omega_1}(A)$ of \Cref{def:Calkin_Small} gives a fiber sequence
		\[
			F(A) \to F(\Ind(A)^{\omega_1}) \to F(\Calk_{\omega_1}(A)).
		\]
		Since $\Ind(A)^{\omega_1}$ admits countable coproducts, the middle term vanishes by an Eilenberg swindle: for $T = \bigoplus_{\N}\colon \Ind(A)^{\omega_1} \to \Ind(A)^{\omega_1}$ we have $\id \oplus T \simeq T$, and additivity gives $F(\id) + F(T) = F(T)$, so $F(\id) = 0$. Hence
		\[
			F(A) \simeq \Omega F(\Calk_{\omega_1}(A)).
		\]
		This is the formula that will be used to extend $F$ to dualizable categories, see \Cref{def:Continuous_Extension}.
	\end{example}

	The most fundamental localizing invariant is nonconnective algebraic K-theory, as made precise by the following universal property.

	\begin{theorem}
		\label[theorem]{thm:K_Universal_Property}
		Nonconnective algebraic K-theory $K\colon \Cat_{\perf} \to \Sp$ is the initial localizing invariant equipped with a natural transformation $(-)^{\simeq} \to \Omega^{\infty}K(-)$.
	\end{theorem}

	\begin{remark}
		\label[remark]{rem:K_Universal_Property_Literature}
		In the literature, \Cref{thm:K_Universal_Property} is usually stated for localizing invariants preserving filtered colimits, see \cite{BGT2013Universal} and also \cite{Barwick_Higher_K_Theory}. The version without this condition is explicitly stated in \cite[page~115, footnote~19]{Ramzi_Thesis}. That footnote does not supply a proof of the additional generality.
	\end{remark}

	\begin{definition}[{\cite{BGT2013Universal}}]
		\label[definition]{def:Mot_Loc}
		The category of \emph{localizing motives} $\mathrm{Mot}_{\loc}$ is the presentable stable category equipped with a localizing invariant
		\[
			U_{\loc}\colon \Cat_{\perf} \to \mathrm{Mot}_{\loc}
		\]
		preserving filtered colimits, characterized as follows: for every presentable stable category $\Ee$, precomposition with $U_{\loc}$ induces an equivalence between colimit-preserving functors $\mathrm{Mot}_{\loc}\to\Ee$ and localizing invariants $\Cat_{\perf}\to\Ee$ preserving filtered colimits. We fix a universe for the small categories in $\Cat_{\perf}$ and use a larger universe for functor categories; these size conventions are suppressed in the notation.
	\end{definition}

	Although $\mathrm{Mot}_{\loc}$ is defined by a universal property, it can be constructed explicitly as a localization of a category of spectral presheaves on a suitable small subcategory of $\Cat_{\perf}$. This is the \emph{finitary} category of motives. It is not asserted to classify all the localizing invariants of \Cref{def:Localizing_Invariant}.

	\begin{theorem}[{\cite[Theorem~9.8]{BGT2013Universal}}]
		\label[theorem]{thm:K_Corepresentable}
		For all $A \in \Cat_{\perf}$ there is a natural equivalence
		\[
			\hom_{\mathrm{Mot}_{\loc}}\big(U_{\loc}(\Sp^{\omega}),U_{\loc}(A)\big) \simeq K(A).
		\]
	\end{theorem}

	% Source: Talk 1 (Marc Hoyois, 13.10.2026)
	Marc formulated the universal property of K-theory as follows.

	\begin{theorem}[{Theorem/definition, cf.\ \cite{BGT2013Universal}}]
		\label[theorem]{thm:K_Unit_Day_Convolution}
		$K$ is the unit for the Day convolution operad structure on $\Fun_{\loc}(\Cat_{\perf},\Sp)$, where localizing invariants are not required to preserve filtered colimits.
	\end{theorem}

	Let us explain the statement. The category $\Fun(\Cat_{\perf},\Sp)$ carries the Day convolution symmetric monoidal structure, induced by the tensor product on $\Cat_{\perf}$ and the smash product on $\Sp$. Its unit is the functor $\Sigma^{\infty}_+\Map_{\Cat_{\perf}}(\Sp^{\omega},-) \simeq \Sigma^{\infty}_+(-)^{\simeq}$. The full subcategory $\Fun_{\loc}(\Cat_{\perf},\Sp)$ of localizing invariants is not obviously closed under Day convolution, but it inherits a structure of an operad: a multimorphism $(F_1,\dots,F_n) \to G$ is a natural transformation $F_1(A_1) \otimes \cdots \otimes F_n(A_n) \to G(A_1 \otimes \cdots \otimes A_n)$. The claim is that $K$ is a unit for this operad: precomposing with the natural transformation $\Sigma^{\infty}_+(-)^{\simeq} \to K$ of \Cref{thm:K_Universal_Property} in one of the inputs induces an equivalence on spaces of multimorphisms whose other inputs and target are localizing.

	% Source: written from the literature (universal property as in Ramzi's thesis); proof drafted with generative AI (Claude)
	\prooffor{thm:K_Unit_Day_Convolution}
	\begin{proof}[Proof sketch]
		By \Cref{thm:K_Universal_Property}, precomposition with $\Sigma^{\infty}_+(-)^{\simeq} \to K$ induces an equivalence
		\[
			\Map(K,H) \simeq \Map\big(\Sigma^{\infty}_+(-)^{\simeq},H\big)
		\]
		for every localizing invariant $H$. Now let $F_2,\dots,F_n$ and $G$ be localizing invariants. A multimorphism $(F_1,F_2,\dots,F_n) \to G$ is the same as a natural transformation from $F_1$ to the functor
		\[
			H(A) := \int_{B_2,\dots,B_n} \hom\big(F_2(B_2) \otimes \cdots \otimes F_n(B_n),\, G(A \otimes B_2 \otimes \cdots \otimes B_n)\big),
		\]
		where $\int$ denotes the end expressing naturality in the remaining variables. The functor $H$ is localizing: for fixed $B_2,\dots,B_n$, the functor $A \mapsto A \otimes B_2 \otimes \cdots \otimes B_n$ preserves short exact sequences, $G$ is localizing, and $\hom(X,-)$ and limits preserve fiber sequences. Applying the above equivalence to $H$ proves the claim for the first input, and by symmetry for every input.
	\end{proof}

	This argument uses the universal property of \Cref{thm:K_Universal_Property} without the assumption that localizing invariants preserve filtered colimits, not the universal property of $\mathrm{Mot}_{\loc}$. Note also that we do not claim that the category of localizing invariants is presentable.

	\begin{theorem}[{Tamme, before Efimov, \cite[Theorem~18]{Tamme_Excision}}]
		\label[theorem]{thm:Tamme_Excision}
		Consider a square $Q$ in $\Cat_{\st}$ of the form
		\[
			\begin{tikzcd}
				A \rar["f"] \dar & B \dar \\
				C \rar["g"'] & D.
			\end{tikzcd}
		\]
		Assume that $\Ind(Q)$ is cartesian, and that $\Ind(f)$ and $\Ind(g)$ are localizations. Then every localizing invariant $F$ sends $Q$ to a cartesian square.
	\end{theorem}

	% Source: proof strategy written from the literature (Efimov, Propositions 2.86 and 8.11)
	We will prove this theorem in \Cref{chap:Localizing_Invariants_Dualizable} for accessible localizing invariants with values in an accessible stable category with $\lambda$-filtered colimits for some $\lambda$, which includes K-theory. The idea is as follows. The square $\Ind(Q)$ is a pullback of compactly generated categories along the localization $\Ind(g)$, so by \Cref{prop:Dualizable_Pullbacks}(2) it is also a pullback in $\Pr^{\dual}$, and $\Ind(f)$ is again a localization with $\ker(\Ind(f)) \simeq \ker(\Ind(g))$. The continuous extension $F^{\cont}$ of \Cref{thm:Efimov_Continuous_Extension} sends the two resulting short exact sequences of dualizable categories to fiber sequences, and $F^{\cont}(\Ind(A)) \simeq F(A)$. Hence $F(Q)$ is cartesian \cite[Proposition~8.11]{Efimov_Large_Categories}. In particular, only $\Ind(g)$ needs to be a localization.

	\begin{warning}
		\label[warning]{warn:Excision_Requires_Ind_Localizations}
		One cannot simply apply $F$ to the common kernel of $f$ and $g$. The hypotheses say that $\Ind(f)$ and $\Ind(g)$ are localizations; they do not say that $f$ and $g$ are Karoubi projections. Their small kernels therefore need not give short exact sequences to which $F$ applies.
	\end{warning}

	\begin{remark}
		\label[remark]{rem:Tamme_Lax_Pullback_Proof}
		Tamme's original proof predates continuous K-theory and applies to localizing invariants with values in any stable category, without accessibility assumptions. It uses the lax pullback $P = B \times_D^{\to} C$ of triples $(b,c,\varphi\colon h(b) \to g(c))$, where $h\colon B \to D$ is the right vertical functor. There is a split exact sequence $C \to P \to B$ \cite[Proposition~10]{Tamme_Excision}, and under the hypotheses of \Cref{thm:Tamme_Excision} the functor $(b,c,\varphi) \mapsto \cofib(\varphi)$ fits into an exact sequence $A \to P \to D$ \cite[Theorem~15]{Tamme_Excision}. Comparing the two sequences gives the claim. See \cite{Land_Tamme_Pullbacks} for a generalization to arbitrary pullbacks of ring spectra.
	\end{remark}

	\begin{example}[{Milnor patching, \cite{Milnor_Intro_K_Theory}}]
		\label[example]{ex:Milnor_Patching}
		Consider a pullback square of $\bbE_1$-rings:
		\[
			\begin{tikzcd}
				A \rar \dar & A' \dar \\
				B \rar & B'
			\end{tikzcd}
		\]
		and assume that $A' \otimes_A A' \iso A'$. Then the induced square
		\[
			\begin{tikzcd}
				\Mod_A \rar \dar & \Mod_{A'} \dar \\
				\Mod_B \rar & \Mod_{B'}
			\end{tikzcd}
		\]
		is also a pullback square, and $\Mod_B \to \Mod_{B'}$ is a localization \cite[Lemma~23 and Theorem~26]{Tamme_Excision}. In particular, we deduce:
	\end{example}

	\begin{theorem}[{\cite[Theorem~28]{Tamme_Excision}; see also \cite{Land_Tamme_Pullbacks}}]
		\label[theorem]{thm:Milnor_Patching_K_Theory}
		In the above situation, there is a pullback square
		\[
			\begin{tikzcd}
				K(A) \rar \dar & K(A') \dar \\
				K(B) \rar & K(B').
			\end{tikzcd}
		\]
	\end{theorem}

	\begin{example}
		\label[example]{ex:Milnor_Square_Of_Rings}
		Given a map of rings $A \to B$ that sends an ideal $I \subseteq A$ isomorphically to the induced ideal $IB \subseteq B$, the induced square
		\[
			\begin{tikzcd}
				A \rar \dar & A/I \dar \\
				B \rar & B/IB
			\end{tikzcd}
		\]
		is a pullback of $\bbE_1$-rings.
	\end{example}
